\section{Probability density functions}

The continuous analogue of a probability mass function is a probability density
function. This sentence is useful, but it can also be misleading if it is read
too quickly. A mass function gives probabilities of individual values:
\[
p_X(x)=P(X=x).
\]
A density function does not usually do that. A density describes how probability
is spread along the real line so that probabilities of intervals can be obtained
by integration.

\begin{definitionbox}
A random variable \(X\) is called \textbf{continuous with density} \(f_X\) if
there is a non-negative function \(f_X:\mathbb{R}\to[0,\infty)\) such that for
every interval \([a,b]\),
\[
P(a\leq X\leq b)=\int_a^b f_X(x)\,dx.
\]
The function \(f_X\) is called a \textbf{probability density function} of \(X\).
\end{definitionbox}

There are two requirements hidden in this definition. First,
\[
f_X(x)\geq 0
\]
for every \(x\), because probabilities cannot be negative. Second, the total
area under the density must be \(1\):
\[
\int_{-\infty}^{\infty} f_X(x)\,dx=1.
\]
This expresses the fact that \(X\) must take some real value with total
probability \(1\).

\begin{remarkbox}
A proposed density is not valid until it is non-negative and has total integral
\(1\). These are the continuous analogues of the two basic properties of a
probability mass function.
\end{remarkbox}

\subsection*{Density is not probability}

The value \(f_X(x)\) is a height. It is not, by itself, the probability that
\(X=x\). The probability of an interval is an area:
\[
P(a\leq X\leq b)=\int_a^b f_X(x)\,dx.
\]
A height becomes a probability only after it is accumulated over a region.

This distinction is essential. If \(f_X(2)=0.4\), this does not mean
\[
P(X=2)=0.4.
\]
It means that near \(x=2\), probability is being distributed with local density
approximately \(0.4\) per unit of length. To turn this into a probability, one
must choose an interval, for instance
\[
P(2\leq X\leq 2.1)=\int_2^{2.1} f_X(x)\,dx.
\]
If the density is nearly constant on that small interval, this probability is
approximately
\[
0.4\cdot 0.1=0.04.
\]

\subsection*{Example: uniform density on \([0,1]\)}

For a point chosen uniformly from \([0,1]\), the density is
\[
f_X(x)=
\begin{cases}
1, & 0\leq x\leq 1,\\
0, & \text{otherwise}.
\end{cases}
\]
The total area is
\[
\int_{-\infty}^{\infty} f_X(x)\,dx=\int_0^1 1\,dx=1.
\]
For an interval \([a,b]\subseteq[0,1]\),
\[
P(a\leq X\leq b)=\int_a^b 1\,dx=b-a.
\]
Thus
\[
P(0.2\leq X\leq 0.5)=\int_{0.2}^{0.5}1\,dx=0.3.
\]
This reproduces the length interpretation from the previous section.

\subsection*{Example: uniform density on \([2,6]\)}

Suppose \(X\) is uniformly distributed on \([2,6]\). The interval has length
\(4\), so the density must have constant height \(1/4\) on that interval:
\[
f_X(x)=
\begin{cases}
\frac14, & 2\leq x\leq 6,\\
0, & \text{otherwise}.
\end{cases}
\]
The total area is
\[
\int_2^6 \frac14\,dx=\frac14(6-2)=1.
\]
For example,
\[
P(3\leq X\leq 5)=\int_3^5 \frac14\,dx=\frac12.
\]
The height \(1/4\) is not a probability. The probability \(1/2\) comes from the
area of a rectangle with width \(2\) and height \(1/4\).

\subsection*{Can a density be larger than one?}

Yes. A density value can be larger than \(1\). This often surprises students
because probabilities cannot be larger than \(1\). But density is not
probability.

For example, consider the uniform distribution on \([0,0.2]\). The interval has
length \(0.2\), so the constant density must be
\[
f_X(x)=5,
\]
for \(0\leq x\leq 0.2\). The total area is still
\[
\int_0^{0.2}5\,dx=5\cdot 0.2=1.
\]
The height is \(5\), but no probability has exceeded \(1\). The probability of a
smaller interval, say \([0,0.05]\), is
\[
P(0\leq X\leq 0.05)=\int_0^{0.05}5\,dx=0.25.
\]

\begin{remarkbox}
The units of a density are ``probability per unit of \(x\)''. If \(x\) is measured
in hours, the density is measured per hour. This is another reason why a density
height is not itself a probability.
\end{remarkbox}

\subsection*{Example: a triangular density}

Let
\[
f_X(x)=
\begin{cases}
2x, & 0\leq x\leq 1,\\
0, & \text{otherwise}.
\end{cases}
\]
This is a valid density because it is non-negative and
\[
\int_0^1 2x\,dx=\left[x^2\right]_0^1=1.
\]
The density is small near \(0\) and larger near \(1\), so values close to \(1\)
are more likely to occur than values close to \(0\), in the sense that intervals
near \(1\) have larger probability than intervals of the same length near \(0\).

For example,
\[
P(0\leq X\leq 0.5)=\int_0^{0.5}2x\,dx=0.25,
\]
whereas
\[
P(0.5\leq X\leq 1)=\int_{0.5}^{1}2x\,dx=0.75.
\]
Both intervals have the same length, but the second has larger probability
because the density is higher there.

\subsection*{What the graph of a density tells us}

A density graph should be read by areas, not by point heights alone. High parts
of the graph indicate regions where probability accumulates quickly. Low parts
indicate regions where probability accumulates slowly. Regions outside the
support, where the density is zero, receive no probability.

When looking at a density, ask:
\begin{itemize}
    \item where is the density non-zero?
    \item is the density constant, increasing, decreasing, or concentrated in a
    particular region?
    \item does the total area under the curve equal \(1\)?
    \item which intervals should have large probability?
    \item which intervals should have small probability?
\end{itemize}

\begin{theorembox}
For a continuous random variable with density \(f_X\), probabilities are computed
by area:
\[
P(X\in A)=\int_A f_X(x)\,dx
\]
for intervals and, more generally, for sets for which the integral is defined.
The value \(f_X(x)\) is a density height, not a point probability.
\end{theorembox}
